\def\tituloconcepto{El qubit como vector}
% Si el capitulo se compila solo (sin main.tex), esta cabecera abre el documento.
% Antes de % Si el capitulo se compila solo (sin main.tex), esta cabecera abre el documento.
% Antes de % Si el capitulo se compila solo (sin main.tex), esta cabecera abre el documento.
% Antes de \input{cabecera} cada capitulo define \tituloconcepto.
\ifdefined\enmain\else
  \documentclass[11pt,a4paper]{article}
  \input{preambulo}
  \begin{document}
  \title{\tituloconcepto}
  \author{\autor}
  \date{\today}
  \maketitle
\fi
 cada capitulo define \tituloconcepto.
\ifdefined\enmain\else
  \documentclass[11pt,a4paper]{article}
  % Preambulo compartido por todos los documentos. Compilar con XeLaTeX/LuaLaTeX
% (o Tectonic). Si se usa Overleaf, elegir el compilador XeLaTeX.
\usepackage[spanish,es-noshorthands]{babel}
\usepackage{amsmath,amssymb}
\usepackage{booktabs,array}
\usepackage[a4paper,margin=2.5cm]{geometry}
\usepackage{xcolor}
\usepackage{listings}
\usepackage[most]{tcolorbox}
\usepackage{quantikz}
\usepackage[hidelinks]{hyperref}

% ---- Autor (editar a gusto) -------------------------------------------------
\newcommand{\autor}{Genesis Silesky}

% ---- Notacion de Dirac y atajos ---------------------------------------------
% quantikz ya define \ket y \bra; DeclareRobustCommand los redefine sin error.
\DeclareRobustCommand{\ket}[1]{\left|#1\right\rangle}
\DeclareRobustCommand{\bra}[1]{\left\langle#1\right|}
\DeclareRobustCommand{\braket}[2]{\left\langle#1\middle|#2\right\rangle}
\newcommand{\C}{\mathbb{C}}
\newcommand{\R}{\mathbb{R}}
\DeclareMathOperator{\Tr}{Tr}

% ---- Cajas de apoyo --------------------------------------------------------
% "concepto": definiciones previas que hacen falta para entender el capitulo.
\newtcolorbox{concepto}[1]{breakable,colback=blue!4,colframe=blue!55!black,
  fonttitle=\bfseries,title={Concepto previo: #1}}
% "nota": aclaraciones y advertencias.
\newtcolorbox{nota}[1]{breakable,colback=orange!6,colframe=orange!70!black,
  fonttitle=\bfseries,title={#1}}

% ---- Estilos de listados ----------------------------------------------------
\lstdefinestyle{py}{language=Python,basicstyle=\ttfamily\footnotesize,
  breaklines=true,postbreak=\mbox{\textcolor{gray}{$\hookrightarrow$}\space},
  numbers=left,numberstyle=\tiny\color{gray},frame=single,
  commentstyle=\color{green!40!black}\itshape,
  keywordstyle=\color{blue!70!black}\bfseries,
  stringstyle=\color{red!60!black},showstringspaces=false,tabsize=4,
  columns=fullflexible,keepspaces=true}
\lstdefinestyle{salida}{basicstyle=\ttfamily\footnotesize,breaklines=true,
  frame=leftline,backgroundcolor=\color{gray!8},columns=fullflexible,
  keepspaces=true}

  \begin{document}
  \title{\tituloconcepto}
  \author{\autor}
  \date{\today}
  \maketitle
\fi
 cada capitulo define \tituloconcepto.
\ifdefined\enmain\else
  \documentclass[11pt,a4paper]{article}
  % Preambulo compartido por todos los documentos. Compilar con XeLaTeX/LuaLaTeX
% (o Tectonic). Si se usa Overleaf, elegir el compilador XeLaTeX.
\usepackage[spanish,es-noshorthands]{babel}
\usepackage{amsmath,amssymb}
\usepackage{booktabs,array}
\usepackage[a4paper,margin=2.5cm]{geometry}
\usepackage{xcolor}
\usepackage{listings}
\usepackage[most]{tcolorbox}
\usepackage{quantikz}
\usepackage[hidelinks]{hyperref}

% ---- Autor (editar a gusto) -------------------------------------------------
\newcommand{\autor}{Genesis Silesky}

% ---- Notacion de Dirac y atajos ---------------------------------------------
% quantikz ya define \ket y \bra; DeclareRobustCommand los redefine sin error.
\DeclareRobustCommand{\ket}[1]{\left|#1\right\rangle}
\DeclareRobustCommand{\bra}[1]{\left\langle#1\right|}
\DeclareRobustCommand{\braket}[2]{\left\langle#1\middle|#2\right\rangle}
\newcommand{\C}{\mathbb{C}}
\newcommand{\R}{\mathbb{R}}
\DeclareMathOperator{\Tr}{Tr}

% ---- Cajas de apoyo --------------------------------------------------------
% "concepto": definiciones previas que hacen falta para entender el capitulo.
\newtcolorbox{concepto}[1]{breakable,colback=blue!4,colframe=blue!55!black,
  fonttitle=\bfseries,title={Concepto previo: #1}}
% "nota": aclaraciones y advertencias.
\newtcolorbox{nota}[1]{breakable,colback=orange!6,colframe=orange!70!black,
  fonttitle=\bfseries,title={#1}}

% ---- Estilos de listados ----------------------------------------------------
\lstdefinestyle{py}{language=Python,basicstyle=\ttfamily\footnotesize,
  breaklines=true,postbreak=\mbox{\textcolor{gray}{$\hookrightarrow$}\space},
  numbers=left,numberstyle=\tiny\color{gray},frame=single,
  commentstyle=\color{green!40!black}\itshape,
  keywordstyle=\color{blue!70!black}\bfseries,
  stringstyle=\color{red!60!black},showstringspaces=false,tabsize=4,
  columns=fullflexible,keepspaces=true}
\lstdefinestyle{salida}{basicstyle=\ttfamily\footnotesize,breaklines=true,
  frame=leftline,backgroundcolor=\color{gray!8},columns=fullflexible,
  keepspaces=true}

  \begin{document}
  \title{\tituloconcepto}
  \author{\autor}
  \date{\today}
  \maketitle
\fi


\section{El qubit como vector}

\subsection{Objetivo}
Entender que un qubit es un vector de dos amplitudes complejas normalizado, y
cómo esas amplitudes determinan las probabilidades de medición.

\subsection{Conceptos previos}
\begin{concepto}{lo que hace falta para este capítulo}
Ver el capítulo de prerrequisitos (sección~\ref{sec:prerreq}): números
complejos (módulo y fase), vectores y norma, notación de Dirac y regla de Born.
Si solo se recuerda una cosa: el \emph{módulo al cuadrado} de una amplitud es
una probabilidad.
\end{concepto}

\subsection{Fundamento teórico}
Un bit clásico vale $0$ o $1$. Un qubit es un vector
\begin{equation}
  \ket{\psi}=\alpha\ket0+\beta\ket1=\binom{\alpha}{\beta},\qquad
  \alpha,\beta\in\C,\quad |\alpha|^2+|\beta|^2=1 .
\end{equation}
Por la regla de Born, al medir se obtiene $0$ con probabilidad $|\alpha|^2$ y
$1$ con probabilidad $|\beta|^2$. La condición de normalización asegura que
las probabilidades sumen 1. Dos observaciones:
\begin{itemize}
  \item \textbf{Fase global:} $e^{i\gamma}\ket\psi$ y $\ket\psi$ son físicamente
    indistinguibles.
  \item \textbf{Fase relativa:} la fase de $\beta$ respecto a $\alpha$ no altera
    las probabilidades de esta medición, pero sí cómo interfiere el qubit con
    otros después (capítulo~05a).
\end{itemize}
Salvo fase global, todo qubit se parametriza como
$\cos\frac\theta2\ket0+e^{i\varphi}\sin\frac\theta2\ket1$ (esfera de Bloch).

\paragraph{Ejemplo del código.} Se toma $\alpha=1/\sqrt3$ y
$\beta=\sqrt{2/3}\,e^{i\pi/4}=(1+i)/\sqrt3$. Entonces
$|\alpha|^2=1/3$ y $|\beta|^2=2/3$, por lo que $P(0)=1/3$, $P(1)=2/3$.

\subsection{Código (completamente comentado)}
\lstinputlisting[style=py]{../../fuentes/01_qubit_vector.py}

\subsection{Explicación celda por celda}
\begin{description}
  \item[Celda 1] Importa \texttt{numpy} (complejos y funciones) y
    \texttt{Statevector}, la clase que guarda las amplitudes.
  \item[Celda 2] Construye $\ket0=(1,0)$ y $\ket1=(0,1)$ con
    \texttt{from\_label}. Son la base canónica.
  \item[Celda 3] Define $\alpha,\beta$ y crea el estado. La suma
    $|\alpha|^2+|\beta|^2$ sale $1.0000000000000002$ y no exactamente $1$:
    es error de redondeo de punto flotante (aritmética de doble precisión).
  \item[Celda 4] \texttt{probabilities\_dict()} aplica la regla de Born. Se
    convierte a \texttt{float} y \texttt{str} para que la impresión sea
    legible (NumPy~2 imprime sus tipos propios).
  \item[Celda 5] Muestra que $(1,1)$ no es válido (norma$^2$ $=2$) y que
    $(1,1)/\sqrt2$ sí lo es.
\end{description}

\subsection{Salida obtenida}
\lstinputlisting[style=salida]{../../fuentes/salidas/01_qubit_vector.txt}
Esta salida es determinista (no hay medición).

\subsection{Interpretación}
Las amplitudes son números complejos, pero las probabilidades salen reales y
suman 1. La fase de $\beta$ ($\pi/4$) no aparece en las probabilidades: es
información que solo se revela mediante interferencia.

\subsection{Ejercicios}
\begin{enumerate}
  \item Cambie la fase de $\beta$ a $\pi$ y verifique que las probabilidades no
    cambian.
  \item Construya un estado con $P(1)=0.9$ y compruébelo.
  \item ¿Es $(\tfrac35,\tfrac45 i)$ un estado válido? Verifíquelo con
    \texttt{is\_valid()}.
\end{enumerate}

\subsection{Errores comunes}
\begin{itemize}
  \item Olvidar normalizar: Qiskit no lo hace automáticamente y
    \texttt{Statevector} podría ser inválido.
  \item Confundir amplitud con probabilidad: la amplitud puede ser negativa o
    compleja; la probabilidad es su módulo al cuadrado.
\end{itemize}

\ifdefined\enmain\else
  \end{document}
\fi

